\scanpage{47}
\mspage{10}
\deleted{(Thus for the \deleted{\unsure{strong}}\inserted{``complete''} Derivator, with \uline{loc.} constant coeff. $\xi$ instead of constant ones, one must suppose that the pr. map $X\longrightarrow\overline X$ is $D$-equiv. $2$-complete.}

To have the converse implication
\[i\in W_{\underline D}\Longrightarrow i_0\in W_{\underline D},\]
one would need to know that the diagram (**), more \deleted{\unsure{part}} precisely the \deleted{\unsure{i\longrightarrow}} abstract \deleted{diagram}\inserted{square} $\sigma$ which gives rise to it, is not only cartesian, but also \uline{cocartesian}. ``Morally'', that it is cartesian seems to me to be the cohomological translation of the fact that the \deleted{diagram}\inserted{square}
\[\diagram{\begin{tikzcd}&U\arrow[dl]&\\U\cap U'&&X\arrow[ul,no head]\arrow[dl]\\&U'\arrow[ul]&\end{tikzcd}}\quad\text{in Cat}\]
is \uline{cocartesian} in a sufficiently strong sense. But it is also cartesian, trivially, and one can expect (under certain conditions) that the diagrams ``transfigured'' by $H^\bullet_{\underline D}(-,\xi)$ are cocartesian\ldots

I feel like looking at $X$ from the outset as over $\underline Q=\Delta^1\times\Delta^1$, rather than over \unsure{the open} $\underline\Psi=\underline Q-\{e\}$ of $\underline Q$, then
\[\underline\Psi_{\underline Q}(X)|\underline\Psi\simeq\Psi_{\underline\Psi}(X)\quad\text{already considered.}\]
\mspage{46}
